%%%%%%%%%%%%%%%%%%%%%%%%%%%%%%%%%%%%%%%%%%%%%%%%%%%%%%%%%%%%%%%%%%%%%%%%%
%                                                                       %
% ustthesis_test.tex: A template file for usage with ustthesis.cls      %
%                                                                       %
%%%%%%%%%%%%%%%%%%%%%%%%%%%%%%%%%%%%%%%%%%%%%%%%%%%%%%%%%%%%%%%%%%%%%%%%%

\documentclass[a4paper]{ustthesis}
\usepackage[square,numbers]{natbib}
\usepackage{xeCJK}
% \usepackage{ctex}
\usepackage[T1]{fontenc}
\usepackage{longtable}
\usepackage{url}
\usepackage{multirow}
\usepackage{graphics}
\usepackage{graphicx}
\usepackage{amsmath}
\usepackage{cases}
\usepackage{colortbl}
\usepackage{xcolor}
\usepackage{times}
\usepackage{array}
\usepackage{pdflscape}
\DeclareMathOperator*{\argmax}{argmax}      % for argmax

% \usepackage{latexsym}
    % Use the "latexsym" package when encountering the following error:
    %   ! LaTeX Error: Command \??? not provided in base LaTeX2e.
% \usepackage{epsf}
    % Use the "epsf" package for including EPS files.

%%%%%%%%%%%%%%%%%%%%%%%%%%%%%%%%%%%%%%%%%%%%%%%%%%%%%%%%%%%%%%%%%%%%%%%%%
%                                                                       %
% Preambles. DO NOT ERASE THEM. Change to suite your particular purpose.%
%                                                                       %
%%%%%%%%%%%%%%%%%%%%%%%%%%%%%%%%%%%%%%%%%%%%%%%%%%%%%%%%%%%%%%%%%%%%%%%%%

\title{Enhancing Reliabilty and Controllability in Generative Models through Domain-informed Constraints}  % Title of the thesis.
\author{He CAO}     % Author of the thesis.
\degree{\PhD}             % Degree for which the thesis is. Options: \AM \MSc \MPhil \PhD
\stage{\Thesis}              % Stage of PhD document; use \Thesis for all other degree. Options: \PQE \Proposal \Thesis
\subject{Individualized Interdisciplinary Program (Artificial Intelligence)} % Subject of the Degree.
\department{Division of Emerging Interdisciplinary Areas}       % Department to which the thesis is submitted.
\advisor{Prof. Yuan YAO}     % Supervisor. Additional co-supervisor can be added using \member
%\acting      % Uncomment for Accting department head
\coadvisor{Prof. Yangqiu SONG}
\depthead{Prof. Huamin QU}     % department head.
\defencedate{2023}{05}{10}     % \defencedate{year}{month}{day}.

% NOTE:
%   According to the sample shown in the guidelines, page number is
%   placed below the bottom margin.  However, if the author prefers
%   the page number to be printed above the bottom margin, please
%   activate the following command.

%\PNumberAboveBottomMargin

\begin{document}

%%%%%%%%%%%%%%%%%%%%%%%%%%%%%%%%%%%%%%%%%%%%%%%%%%%%%%%%%%%%%%%%%%%%%%%%%
%                                                                       %
% Now the actual Thesis. The order of output MUST be followed:          %
%                                                                       %
%    1) TITLEPAGE                                                       %
%                                                                       %
% The \maketitle command generates the Title page as well as the        %
% Signature page.                                                       %
%                                                                       %
%%%%%%%%%%%%%%%%%%%%%%%%%%%%%%%%%%%%%%%%%%%%%%%%%%%%%%%%%%%%%%%%%%%%%%%%%

\maketitle

%%%%%%%%%%%%%%%%%%%%%%%%%%%%%%%%%%%%%%%%%%%%%%%%%%%%%%%%%%%%%%%%%%%%%%%%%
%                                                                       %
%     2) DEDICATION (Optional)                                          %
%                                                                       %
% The \dedication and \enddedication commands are optional. If          %
% specified it generates a page for dedication.                         %
%
%%%%%%%%%%%%%%%%%%%%%%%%%%%%%%%%%%%%%%%%%%%%%%%%%%%%%%%%%%%%%%%%%%%%%%%%%

\dedication
% This is an optional section.
Here is the dedication.
\enddedication
\newpage

%%%%%%%%%%%%%%%%%%%%%%%%%%%%%%%%%%%%%%%%%%%%%%%%%%%%%%%%%%%%%%%%%%%%%%%%%
%                                                                       %
%     3) ACKNOWLEDGMENTS                                                %
%                                                                       %
% \acknowledgments and \endacknowledgments defines the                  %
% Acknowledgments of the author of the Thesis.                          %
%                                                                       %
%%%%%%%%%%%%%%%%%%%%%%%%%%%%%%%%%%%%%%%%%%%%%%%%%%%%%%%%%%%%%%%%%%%%%%%%%

\acknowledgments
Thanks your supervisor(s) and the defense committee. 

Thanks everyone help in your thesis research.

Thanks everyone support you throughout the way.

Thanks your loved ones.

\endacknowledgments
\newpage

%%%%%%%%%%%%%%%%%%%%%%%%%%%%%%%%%%%%%%%%%%%%%%%%%%%%%%%%%%%%%%%%%%%%%%%%%
%                                                                       %
%     4) TABLE OF CONTENTS                                              %
%                                                                       %
%%%%%%%%%%%%%%%%%%%%%%%%%%%%%%%%%%%%%%%%%%%%%%%%%%%%%%%%%%%%%%%%%%%%%%%%%

\tableofcontents

%%%%%%%%%%%%%%%%%%%%%%%%%%%%%%%%%%%%%%%%%%%%%%%%%%%%%%%%%%%%%%%%%%%%%%%%%
%                                                                       %
%     5) LIST OF FIGURES (If Any)                                       %
%                                                                       %
%%%%%%%%%%%%%%%%%%%%%%%%%%%%%%%%%%%%%%%%%%%%%%%%%%%%%%%%%%%%%%%%%%%%%%%%%

\listoffigures

%%%%%%%%%%%%%%%%%%%%%%%%%%%%%%%%%%%%%%%%%%%%%%%%%%%%%%%%%%%%%%%%%%%%%%%%%
%                                                                       %
%     6) LIST OF TABLES (If Any)
%                                                                       %
%%%%%%%%%%%%%%%%%%%%%%%%%%%%%%%%%%%%%%%%%%%%%%%%%%%%%%%%%%%%%%%%%%%%%%%%%

\listoftables

%%%%%%%%%%%%%%%%%%%%%%%%%%%%%%%%%%%%%%%%%%%%%%%%%%%%%%%%%%%%%%%%%%%%%%%%%
%                                                                       %
%     7) ABSTRACT                                                       %
%                                                                       %
% \abstract and \endabstract are used to define a short Abstract for    %
% the Thesis.                                                           %
%                                                                       %
%%%%%%%%%%%%%%%%%%%%%%%%%%%%%%%%%%%%%%%%%%%%%%%%%%%%%%%%%%%%%%%%%%%%%%%%%

\abstract
Abstract text, 300 words or less
\endabstract

%%%%%%%%%%%%%%%%%%%%%%%%%%%%%%%%%%%%%%%%%%%%%%%%%%%%%%%%%%%%%%%%%%%%%%%%%
%                                                                       %
%     8) The Actual Contents                                            %
%                                                                       %
% The command \chapters MUST BE USED to ensure that the entire content  %
% of the Thesis is double-spaced (in version 1.0).                      %
%                                                                       %
% However, in version 2.0, \chapters will be automatically added in     %
% the beginning of the first chapter.                                   %
%                                                                       %
%%%%%%%%%%%%%%%%%%%%%%%%%%%%%%%%%%%%%%%%%%%%%%%%%%%%%%%%%%%%%%%%%%%%%%%%%

%%\chapters         % Not necessary with ustthesis.cls (v2.0).

%%%%%%%%%%%%%%%%%%%%%%%%%%%%%%%%%%%%%%%%%%%%%%%%%%%%%%%%%%%%%%%%%%%%%%%%%
%                                                                       %
% Each chapter is defined via the \chapter command. The usual sectional %
% commands of LaTeX are also available.                                 %
%                                                                       %
%%%%%%%%%%%%%%%%%%%%%%%%%%%%%%%%%%%%%%%%%%%%%%%%%%%%%%%%%%%%%%%%%%%%%%%%%

\chapter{Introduction}
\label{chp_intro}
\section{Motivation}


\section{Thesis Organization and Contributions}

% Sample page: A page of illustrations with legends.
% \begin{figure}[htb]
% \centering
% \includegraphics[width=0.5\columnwidth]{figure/fig_HKUST.png}
% \caption{Standard University Combination Logo - English}
% \label{fig_HKUST}
% \end{figure}

% \begin{table}[htb]
% \caption{Symbolic meaning of the HKUST emblem.}
% \centering

% \begin{tabular}{|c|l|}
% \hline Parts & Meaning \\
% \hline
% \multirow{4}{*}{As a whole}& A scholar holding up an open book \\ 
% & A transmission tower \\
% & - representing engineering and technology \\
% & - or communication and management \\ \hline
% \multirow{2}{*}{Top circle}& Golden head of wisdom of the scholar\\
% & A sun radiating gold (very traditional Chinese color)\\ \hline
% \multirow{5}{*}{Wave beneath the circle} & The open book of knowledge\\
% & An ocean representing Clear Water Bay, Hong Kong \\
% & - where the University locates\\ 
% & Upsilon \\ 
% & Top part of Tau \\\hline
% \multirow{2}{*}{Curves on both side}  & Arms of a scholar\\
% & Sides of Psi \\ \hline
% \multirow{4}{*}{Bottom center}  & Body of the scholar\\
% & Flask representing science\\
% & Center part of Psi \\
% & Bottom part of Tau \\ \hline
% \end{tabular}
% \label{tbl_wmt17-segrr}
% \end{table}
\newpage

\chapter{Background}
\label{chp_background}
% Here is a version of the guidelines on thesis preparation with minor edits. Please refer to the original document \citep{guidelines} for the most update and accurate information about thesis preparation, especially about the thesis submission protocol. 
% \section{Introduction}

\citep{guidelines}
\newpage

%%%%%%%%%%%%%%%%%%%%%%%%%%%%%%%%%%%%%%%%%%%%%%%%%%%%%%%%%%%%%%%%%%%%%%%%%
%                                                                       %
%      9) BIBLIOGRAPHY                                                  %
%                                                                       %
% This example uses bibtex to generate the required Bibliography. Refer %
% to the % the file ustthesis_test.bib for the entries of the           %
% Bibliography. Note that only the cited entries are printed.           %
%                                                                       %
% If BibTeX is not used to typeset the bibliography, replace the        %
% following line with the \begin{thebibliography} and \end{bibliography}%
% commands (the "thebibliography" environment) to process the           %
% Bibliography.                                                         %
%                                                                       %
%%%%%%%%%%%%%%%%%%%%%%%%%%%%%%%%%%%%%%%%%%%%%%%%%%%%%%%%%%%%%%%%%%%%%%%%%

%%%%%%%%%%%%%%%%%%%%%%%%%%%%%%%%%%%%%%%%%%%%%%%%%%%%%%%%%%%%%%%%%%%%%%%%%
%                                                                       %
% The recommended bibliography style is the IEEE bibliography style.    %
% "ustbib" defines the IEEE bibliography standard with the added        %
% ability of sorting the items by name of author.                       %
%                                                                       %
% If you are not using BibTeX to process your Bibliography, comment out %
% the following line.                                                   %
%                                                                       %
%%%%%%%%%%%%%%%%%%%%%%%%%%%%%%%%%%%%%%%%%%%%%%%%%%%%%%%%%%%%%%%%%%%%%%%%%

\addcontentsline{toc}{chapter}{Bibliography}
\bibliographystyle{IEEEtran}
\bibliography{references}
\newpage

%%%%%%%%%%%%%%%%%%%%%%%%%%%%%%%%%%%%%%%%%%%%%%%%%%%%%%%%%%%%%%%%%%%%%%%%%
%                                                                       %
%     10) APPENDIX (If Any)                                              %
%                                                                       %
% \appendix command marks the beginning of the APPENDIX part of the     %
% Thesis. The usual \chapter command is used for the different chapters %
% of the Appendix.                                                      %
%                                                                       %
%%%%%%%%%%%%%%%%%%%%%%%%%%%%%%%%%%%%%%%%%%%%%%%%%%%%%%%%%%%%%%%%%%%%%%%%%
\appendix
\chapter{Tips on HKUST thesis preparation}
\label{chp_tips}
\section{Sideway page}
Here is an example to have a page turned side way in the PDF. It is useful for having wide figures and tables in the thesis. 

\begin{landscape} 
\begin{table}
\caption{Pearson's
        $\rho $        correlation of participated metrics with the WMT 2016 official average
        direct assessment human judgments on newstest 10k hybrid super-sampled
        systems at system level. Correlations of metrics not significantly
        outperformed by any other for that language pair are highlighted in
        bold.}
\centering

\begin{tabular}{m{0.36\textwidth}cccccccc|cccccccc}
\hline & \multicolumn{8}{c|}{\textbf{into-English}} & \multicolumn{8}{c}{\textbf{out-of-English}} \\\textbf{Metric} & \textbf{cs} & \textbf{de} & \textbf{fi} & \textbf{lv} & \textbf{ru} & \textbf{tr} & \textbf{zh} & \textbf{avg.} & \textbf{cs} & \textbf{de} & \textbf{fi} & \textbf{lv} & \textbf{ru} & \textbf{tr} & \textbf{zh} & \textbf{avg.} \\ \hline
Blend  & .963 & \textbf{.969} & .956 & .976 & \textbf{.957} & .981 & .890 & .956 & -- & -- & -- & -- & .950 & -- & -- & -- \\ \hline
BEER & .966 & .952 & .954 & .974 & .930 & .969 & .897 & .949 & .963 & .829 & .975 & .923 & .942 & .968 & .906 & .929 \\ \hline
UHH\_TSKM & .990 & .929 & .918 & \textbf{.986} & .908 & .982  & .896  & .944  & --  & --  & --  & --  & --  & --  & --  & --  \\ \hline
CDER & .983  & .922  & .925  & .981  & .916  & .970  & \textbf{.898 } & .942  & .958  & .803  & .962  & .911  & .922  & .948  & \textbf{.975 } & .926  \\ \hline
TreeAggreg & .977  & .913  & \textbf{.975} & .983  & .912  & \textbf{.983 } & .854  & .942  & .942  & .765  & .963  & .915  & .919  & .971  & .933  & .915  \\ \hline
chrF++ & .935  & .957  & .924  & .970  & .938  & .957  & .876  & .937  & .966  & .835  & .977  & .944  & .942  & .975  & .968  & .944  \\ \hline
chrF & .933  & .960  & .935  & .965  & .946  & .941  & .855  & .934  & .968  & .845  & \textbf{.979 } & .945  & .947  & \textbf{.980 } & .969  & .947  \\ \hline
NIST & \textbf{.994} & .917  & .928  & .957  & .904  & .969  & .831  & .929  & .954  & .761  & .957  & .914  & .917  & .976  & .968  & .921  \\ \hline
TER & .983  & .899  & .950  & .967  & .905  & .951  & .837  & .927  & .951  & .790  & .959 & .888 & .930 & .958 & .965 & .920 \\ \hline
BLEU & .964  & .914  & .906  & .974  & .907  & .969  & .852  & .927  & .945  & .793  & .919  & .839  & .893  & .916  & .969  & .896  \\ \hline
MEANT-dvw.cos.wmax.n2.m.r8.$\alpha $1.0.$\beta $0.1 & .921  & .942  & .939  & .967  & .955  & .931  & .836  & .927  & --  & .844  & --  & --  & --  & --  & .944  & --  \\ \hline
CharacTER & .963  & .965  & .944  & .927  & .948  & .946  & .740  & .919  & \textbf{.973 } & \textbf{.893 } & .970  & .892  & .929  & .961  & .914  & .933 \\ \hline
WER & .981  & .889  & .946  & .965  & .900  & .922  & .828  & .919  & .949  & .797  & .959  & .884  & .931  & .947  & .951  & .917  \\ \hline
PER & .967  & .920  & .892  & .958  & .904  & .898  & .866  & .915  & .960  & .680  & .939  & .817  & .876  & .955  & .893  & .874  \\ \hline
MEANT-dvw.cos.wmax.n2.m.r8.$\alpha $1.0.$\beta $0.0 & .896  & .928  & .931  & .960  & .952  & .895  & .799  & .909  & .968  & .753  & .970  & \textbf{.947 } & \textbf{.955 } & .980  & .931  & .929  \\ \hline
AutoDA & .440  & .951  & .922  & .970  & .902  & .914  & .734  & .833  & .967  & .602  & .879  & .731  & .850  & .586  & .968  & .797  \\ \hline
\end{tabular}
\label{tbl_wmt17-sys}
\end{table}
\end{landscape} 


%%%%%%%%%%%%%%%%%%%%%%%%%%%%%%%%%%%%%%%%%%%%%%%%%%%%%%%%%%%%%%%%%%%%%%%%%
%                                                                       %
%     11) BIOGRAPHY (optional)                                          %
%                                                                       %
% \biography and \endbiography are used to define the optional          %
% Biography of the author of the Thesis.                                %
%                                                                       %
%%%%%%%%%%%%%%%%%%%%%%%%%%%%%%%%%%%%%%%%%%%%%%%%%%%%%%%%%%%%%%%%%%%%%%%%%

\end{document}
